\subsection*{14. Assessment Needs a Domain Model}

To measure expertise, first describe the professional domain. What work actually matters? Which mistakes cause meaningful harm? Which decisions distinguish a novice from a competent practitioner? Who is qualified to judge them? Can the tasks be safely represented in a simulation or must they be observed in the field?

In technical education, consult experienced practitioners and inspect real but appropriately anonymized failure cases. In healthcare, laboratory work, regulated finance, or safety-critical engineering, additional professional and legal requirements apply. The curriculum's convenience must not define the domain. A school that teaches only what is easy to examine will eventually certify proficiency in its own exam rather than the profession.

\subsection*{15. Design a Rubric for Consequential Decisions}

A useful rubric describes several dimensions of competence and recognizable performance levels. For a production-oriented AI evaluation task, 1 could assess problem framing, validity of evidence, implementation reliability, risk awareness, independence, and communication. Each dimension needs anchors: what does unacceptable, emerging, competent, and advanced work actually look like?

The scoring rules should protect essential competencies from compensation by irrelevant strengths. An attractive user interface should not compensate for a fundamentally invalid evaluation. An eloquent explanation should not compensate for a system that leaks confidential data. Where certification is high-stakes, a policy might require minimum performance on every critical dimension rather than a single average that hides fatal gaps.

The threshold must be justified through practitioner judgment, pilot work, and review of the consequences of misclassification. ``80\% equals mastery'' is not a scientific fact independent of test difficulty and intended use.

\subsection*{16. Avoid Benchmark Contamination and Rehearsal}

If learners rehearse the exact final case repeatedly, successful performance may reflect memory of the case rather than broader competence. If an AI tutor has access to protected test items, it may leak solutions. If instructors coach the final project too closely, the finished artifact may reflect their judgment rather than the learner's.

Use fresh or equivalent tasks with the same underlying demands, rotate problem contexts, include plausible distractions, and require independent explanations. Do not declare 2 tasks equivalent merely because both concern the same topic. Check difficulty using expert review and pilot data. Preserve enough stable items or calibrated anchors to compare cohorts without allowing wholesale memorization of the examination.

\subsection*{17. The One-Score Fallacy}

A single score is convenient for ranking, marketing, and certification. It can be a poor model of professional capability. 1 learner may excel at software reliability and struggle with statistical inference; another may do the reverse. Their aggregate scores could be identical while their suitable roles and development needs differ substantially.

For high-value technical programs, consider a \textbf{capability profile} that records essential dimensions, entry evidence, independent final evidence, delayed checks, and the precise role of assistance. Keep the report interpretable to outsiders. A 20-page scorecard nobody understands is not inherently superior to a concise, meaningful summary. The point is to reveal consequential strengths and gaps rather than disguise them with a decorative grade.

\subsection*{18. Calibrate Raters Rather Than Assuming Expertise Produces Agreement}

Two strong engineers may disagree about what counts as a good design. Some disagreement can reveal legitimate alternatives; some reflects ambiguous criteria. Before high-stakes scoring, have reviewers independently rate shared anchor artifacts, discuss reasons, refine descriptors, and retry. Record the rate and nature of serious disagreements.

A simple raw agreement percentage can be useful, but it does not account for chance agreement, class imbalance, or ordinal severity. Cohen's kappa or intraclass correlation may help in appropriate designs, but a sophisticated statistic cannot rescue a badly defined rubric. Analyze \emph{what} reviewers disagree about. If disagreement clusters around ``professional judgment,'' the construct may need clearer criteria or more observations rather than forcing a mechanically precise score.

\subsection*{19. Sample Size Is Part of the Validity of a Claim}

Suppose 9 of 10 participants pass an assessment. Saying ``90\% of our students mastered the skill'' hides considerable uncertainty. The observed proportion is 0.9 in a sample of 10. A different next group could yield a substantially different rate. If the 10 were selected by the founder from unusually prepared contacts, generalization is narrower still.

For simple binomial proportions, uncertainty intervals can be computed, but they rely on assumptions about the sampling and independent observations. In small pilots, report \textbf{counts, denominators, baseline conditions, assessment details, and uncertainty} rather than dressing a tiny group in enterprise-sized percentages. No confidence interval eliminates biased recruitment or invalid assessment.

\subsection*{20. Delayed Performance and Transfer Require Separate Evidence}

An end-of-course performance may reflect intensive recent practice. To assess durability, revisit essential capabilities after a meaningful delay. To assess transfer, change task context or conditions while preserving the governing skill. These are distinct tests. Someone may remember a procedure after 6 weeks yet fail when the problem arrives in a different format.

Define the target retention interval and breadth of transfer in advance. A short programming course need not promise lifelong mastery; a safety-related skill may require periodic recertification. Report follow-up participation because graduates who return for voluntary assessment may not represent all graduates. \textbf{Long-term competence is not proven by 1 successful final afternoon.}
